%! TeX root = ../main.tex%
\section{Исследование высших гармоник сигнала}

В данном разделе исследуется спектральный состав реального сигнала генератора Advantex SG8 в расширенном частотном диапазоне с целью обнаружения и фиксации высших гармонических составляющих (2-й и 3-й гармоник).
\subsection{Методика поиска гармоник и результаты измерений}

В ходе предыдущих измерений точная частота основного колебания генератора, зафиксированная маркером анализатора спектра, составила $f_1 = 999{,}9789\text{ МГц} \approx 1{,}0\text{ ГГц}$.

Поскольку автоматическая подстройка AutoTune сужает полосу обзора вокруг несущей и принципиально не позволяет зафиксировать высшие гармоники, поиск гармонических составляющих проводился вручную путём аналитического расчёта их частот и ступенчатой перестройки центральной частоты прибора. 

Частота второй гармоники определяется удвоением частоты первой гармоники:
\begin{equation}
    f_2 = 2 \cdot f_1 = 2 \cdot 999{,}9789\text{ МГц} = 1999{,}9578\text{ МГц} \approx 2{,}0\text{ ГГц}.
\end{equation}

Частота третьей гармоники определяется умножением частоты основного тона на три:
\begin{equation}
    f_3 = 3 \cdot f_1 = 3 \cdot 999{,}9789\text{ МГц} = 2999{,}9367\text{ МГц} \approx 3{,}0\text{ ГГц}.
\end{equation}

Для каждой из рассчитанных гармоник на анализаторе спектра вручную задавались соответствующие значения центральной частоты настройки $\text{Center Freq}$, а параметры фильтрации $\text{RBW}$ и $\text{VBW}$ устанавливались в автоматический режим. Полученные спектрограммы 2-й и 3-й гармоник приведены на рис.~\ref{fig:second_harmonic} и рис.~\ref{fig:third_harmonic}.

\begin{figure}[H]
    \centering
    \begin{minipage}[b]{0.48\textwidth}
        \centering
        \includegraphics[width=\textwidth]{fig10.png}
        \caption{Спектрограмма 2-й гармоники ($f_2 \approx 2{,}0\text{ ГГц}$)}
        \label{fig:second_harmonic}
    \end{minipage}
    \hfill
    \begin{minipage}[b]{0.48\textwidth}
        \centering
        \includegraphics[width=\textwidth]{fig11.png}
        \caption{Спектрограмма 3-й гармоники ($f_3 \approx 3{,}0\text{ ГГц}$)}
        \label{fig:third_harmonic}
    \end{minipage}
\end{figure}


Теоретический спектр идеального моногармонического колебания (синусоиды) строго локализован на единственной частоте и не содержит дополнительных составляющих. Однако в реальном физическом эксперименте на частотах порядка $2\text{ ГГц}$ и $3\text{ ГГц}$ отчетливо фиксируются спектральные пики. Их появление непосредственно связано с физикой работы схемотехники генератора.

В основе любого высокочастотного источника сигнала лежит автогенератор, ключевыми узлами которого являются колебательная система и нелинейный активный элемент (транзистор), охваченный цепью положительной обратной связи. Из-за нелинейного элемента форма генерируемого тока неизбежно отклоняется от идеальной синусоиды, что по теореме Фурье порождает дискретный набор гармоник на кратных частотах $n \cdot f_1$.

